\section{Methodology}

\subsection{Data Description}

\subsubsection{Bus Dataset}
The bus dataset contains 53,155 GPS records with the following attributes:
\begin{itemize}
    \item \textbf{Spatial}: Latitude, longitude
    \item \textbf{Kinematic}: Speed, bearing (direction)
    \item \textbf{Temporal}: Timestamps
    \item \textbf{Operational}: Door state (open/closed)
\end{itemize}

\subsubsection{Passenger Dataset}
The passenger dataset contains 50,601 smartphone records with:
\begin{itemize}
    \item \textbf{GPS}: Latitude, longitude, speed, bearing
    \item \textbf{Accelerometer}: X, Y, Z acceleration (accx, accy, accz)
    \item \textbf{Gyroscope}: X, Y, Z rotation (rotx, roty, rotz)
    \item \textbf{Magnetometer}: X, Y, Z magnetic field (magx, magy, magz)
    \item \textbf{Activity}: Motion classification (stationary, walking, running, automotive, cycling, unknown)
    \item \textbf{Proximity}: RSSI and proximity sensor readings (rssiA-C, rssi1-2, proxA-C, prox1-2)
    \item \textbf{Web coordinates}: x\_web, y\_web (projected coordinates)
\end{itemize}

\subsection{Feature Engineering Pipeline}

The feature engineering process transforms raw GPS traces into predictive features through multiple stages:

\subsubsection{Stage 1: Temporal Alignment}
Bus and passenger data are temporally aligned using a tolerance window (default: 60 seconds) to match passenger records with corresponding bus positions.

\subsubsection{Stage 2: Kinematic Features}
\paragraph{Speed Features}
\begin{itemize}
    \item Instantaneous speed from GPS
    \item Speed difference between passenger and nearest bus
    \item Speed ratio (passenger/bus)
\end{itemize}

\paragraph{Acceleration Features}
Computed using consecutive GPS points:
\begin{equation}
    a = \frac{\Delta v}{\Delta t}
\end{equation}
where $\Delta v$ is velocity change and $\Delta t$ is time interval.

\paragraph{Bearing Features}
Direction of movement calculated from consecutive coordinates:
\begin{equation}
    \theta = \arctan2(\Delta \text{lon}, \Delta \text{lat})
\end{equation}
Bearing difference between passenger and bus indicates alignment/divergence.

\subsubsection{Stage 3: Proximity Features}
\paragraph{Distance to Nearest Bus}
Haversine distance between passenger and closest bus:
\begin{equation}
    d = 2r \arcsin\left(\sqrt{\sin^2\left(\frac{\Delta\phi}{2}\right) + \cos(\phi_1)\cos(\phi_2)\sin^2\left(\frac{\Delta\lambda}{2}\right)}\right)
\end{equation}
where $\phi$ is latitude, $\lambda$ is longitude, and $r$ is Earth's radius.

\paragraph{Distance to Bus Stops}
Minimum distance to predefined bus stop locations from OpenStreetMap data.

\subsubsection{Stage 4: Sensor Features}
Raw sensor data from smartphones provide additional context:
\begin{itemize}
    \item Accelerometer magnitude: $\|\vec{a}\| = \sqrt{a_x^2 + a_y^2 + a_z^2}$
    \item Activity confidence scores
    \item Proximity and RSSI values
\end{itemize}

\subsection{Dimensionality Reduction and Clustering}

\subsubsection{PCA (Principal Component Analysis)}
With 40+ engineered features, dimensionality reduction prevents overfitting and improves clustering performance. PCA transforms features to 4 principal components capturing maximum variance:

\begin{equation}
    \mathbf{Z} = \mathbf{X} \mathbf{W}
\end{equation}

where $\mathbf{X}$ is the feature matrix and $\mathbf{W}$ is the projection matrix.

\subsubsection{HDBSCAN Clustering}
HDBSCAN operates on the PCA-reduced feature space:

\begin{enumerate}
    \item \textbf{Core distance computation}: For each point, find the distance to its $k$-th nearest neighbor (min\_samples parameter)
    \item \textbf{Mutual reachability graph}: Connect points based on mutual reachability distance
    \item \textbf{Hierarchy construction}: Build minimum spanning tree of the graph
    \item \textbf{Cluster extraction}: Extract stable clusters across density levels
    \item \textbf{Noise labeling}: Points not belonging to any cluster are labeled as noise (label = -1)
\end{enumerate}

Key parameters:
\begin{itemize}
    \item \texttt{min\_cluster\_size}: Minimum points to form a cluster (default: 1000)
    \item \texttt{min\_samples}: Points in neighborhood for core point (default: 5)
\end{itemize}

\subsection{Label Mapping}

HDBSCAN produces cluster labels $\{-1, 0, 1, 2, ...\}$ where:
\begin{itemize}
    \item $-1$ = Noise (outliers)
    \item $0, 1, 2, ...$ = Cluster IDs
\end{itemize}

These are mapped to binary classification:
\begin{itemize}
    \item \textbf{Label 1 (IN)}: Passengers in dominant cluster (typically cluster 1)
    \item \textbf{Label 0 (OUT)}: Noise and other clusters
\end{itemize}

The assumption is that passengers inside the bus form a dense, cohesive cluster in the feature space due to similar kinematic and proximity characteristics.

\subsection{Pipeline Architecture}

The complete pipeline follows scikit-learn's transformer API:

\begin{lstlisting}[style=python, caption=Pipeline Structure]
from sklearn.pipeline import Pipeline

pipeline = Pipeline([
    ('speed', SpeedTransformer()),
    ('acceleration', AccelerationTransformer()),
    ('bearing', BearingTransformer()),
    ('proximity', ProximityTransformer()),
    ('pca_hdbscan', PCAHDBSCANTransformer()),
    ('label_mapper', LabelMapperTransformer())
])
\end{lstlisting}

This modular design enables:
\begin{itemize}
    \item Independent testing of each component
    \item Easy hyperparameter tuning
    \item Reproducible preprocessing
    \item Seamless integration with scikit-learn tools
\end{itemize}
